
%\usepackage{comment}





\section{Fractional Algorithm for Strengthened LP}%\textsf{Stronger-LP}}
\label{sec:stronger}


The integrality gap example shows that LP \eqref{eq:LP} is not sufficient for obtaining a randomized $\textnormal{polylog}(\ell)$-competitive algorithm for general non-linear paging. Instead, we describe a stronger version of our LP \eqref{eq:LP}, %
in which we require to remove from each infeasible set $S$ a set of pages $S'$, so that the complement of $S'$ in $S$, $S \setminus S'$, will be feasible in the cache (i.e., $ f\left(S \setminus S'\right) \leq k$).  
%As we do not know a priori which such set $S'$ is evicted by the integral optimum, we only demand that evicting a minimum number of pages from $S$ whose complement size fits in cache. 
Formally, for a set $S \subseteq \cP$, let $\cC(S) = \left\{S' \subseteq S \text{ s.t. } f\left(S \setminus S'\right) \leq k\right\}$ be the collection of subsets of $S$ such that $ f\left(S \setminus S'\right) \leq k$ and define
\begin{equation}
	\label{eq:q(S)}
	q(S) = \min_{S' \in \cC(S)} \left|S'\right|
\end{equation}
as the number of pages needed to be evicted from $S$ at any point in time. 
The LP is given below using similar notation to \eqref{eq:LP}.


\begin{equation}
	\label{eq:LPS}
	\begin{aligned}
		\textsf{Stronger-LP}:~~~~~~~~	& \min \sum_{p \in \cP~} \sum_{j \in [n_p]}  x_p(j) \cdot c(p)\\
		& ~~~~~~~~\text{s.t. }\\
		& \sum_{p \in S-p_t} x_p(r(p,t)) \geq q(S), 	~~~~~~~~~~~~~~\forall t \in T~\forall S  \in \cS(t) \\ 
		& x_p(j) \geq 0 ~~~~~~~~~~~~~~~~~~~~~~~~~~~~~~~~~~~	\forall p \in \cP~\forall j \in [n_p]
	\end{aligned}
\end{equation} For example, in classic paging, for a subset of pages $S \subseteq \cP$ it holds that $q(S) = |S|-k$, and hence the above constraints require removing at least $|S|-k$ pages from cache within $S$, coinciding with the LP used for solving (weighted) paging \cite{bansal2012primal}. % at least for classic paging, the integrality gap of the LP does not exceed $O(\log k)$. 
Clearly, LP \eqref{eq:LPS} is a relaxation of non-linear paging, since an integral solution, for any infeasible set $S$, must evict at least $q(S)$ pages from $S$. 








In this section, we give a fractional $O\left(\log \mu\right)$-competitive algorithm for solving \textsf{Stronger-LP} given in \eqref{eq:LPS}, which yields the proof of \Cref{thm:MU}. The dual LP of \eqref{eq:LPS} is as follows. 

\begin{equation}
	\label{eq:dualS}
	\begin{aligned}
		%\textsf{Dual-LP}:
		~~~~~~~~	& \max~~~ \sum_{t \in T~} \sum_{S  \in \cS(t)} y_t(S) \cdot q(S)\\
		& ~~~~~~~~\text{s.t. }\\
		& \sum_{t \in I(p,j)~} \sum_{S \in \cS(t)\big|p \in S-p_t} y_t(S) \leq c(p), 	~~~~~~~~~~~~~~\forall p \in \cP~\forall j \in [n_p]\\ 
		%	& y_t(S) \geq 0 ~~~~~~~~~~~~~~~~~~~~~~~~~~~~~	\forall t \in T(I) ~\forall S \in \bbc. 
	\end{aligned}
\end{equation} 


Let $\textsf{Primal-LP}$ and $\textsf{Dual-LP}$  denote the optimum values of optimal solutions for the primal and the dual programs, respectively.  %, and by $\OPT$ the offline (integral) optimum.  
Our fractional algorithm initializes the (infeasible) primal solution and the (feasible) dual solutions as vectors of zeros $\bar{0}$. Upon the arrival of the requested page $p_t$ at time $t$ we do the following process until $x$ satisfies all constraints of the LP \eqref{eq:LPS} up to time $t$. 


\begin{algorithm}[h]
	\caption{$\textsf{Fractional}$}
	\label{alg:fractionalS}
	
	%	Initialize the cache $\cG_0 \leftarrow \emptyset$.
	Initialize (infeasible) primal solutions $x,z \leftarrow \bar{0}$ 
	
	Initialize (feasible) dual solution $y \leftarrow \bar{0}$.  
	
	
	\For{$t \in T$}{
		
		\While{$x$ is not feasible for $t$\label{step:Whilefeasible}}{
			
			%	Let $\textsf{Frac}_t = \left\{p \in \cP~|~r(p,t) \geq 1 \textnormal{ and } x_p(r(p,t)) < 1\right\}$ be pages not fully evicted.   
			
			Find $S \subseteq \cS(t)$ %of maximum cardinality 
			such that $\sum_{p \in S-p_t} x_p(r(p,t)) < q(S)$ and $x_p(r(p,t))<1~\forall p \in S-p_t $.\label{step:fQS}
			
			
			\While{$\sum_{p \in S-p_t} x_p(r(p,t)) < q(S)$ \textnormal{and} $x_p(r(p,t))<1~\forall p \in S-p_t $\label{step:InWhileS}}{
				
				Increase $y_t(S)$ continuously.\label{step:fyS}
				
				
				\ForAll{$p \in S-p_t$}{
					
					increase $x_p(r(p,t))$ according to\label{step:fxS}
					%$$\frac{\ln (k+1) \cdot }{c(p)}$$
					$$x_p(r(p,t)) \leftarrow \frac{1}{\mu} \cdot \left(     \exp \left(  \frac{\ln (n+1)}{c(p)} \cdot Y_p(r(p,t))\right)-1\right).$$
					
					
					
					
%					\If{$x_p(r(p,t))-\frac{z_p(r(p,t))}{2}\geq \frac{1}{4 \cdot N \cdot n}$}{
%						
%						$z_p(r(p,t)) \leftarrow 2 \cdot x_p(r(p,t))$.\label{step:zS}
%						
%					}
%					
%					
%					\If{$x_p(r(p,t)) \geq \frac{1}{2}$}{
%						
%						$z_p(r(p,t)),x_p(r(p,t))  \leftarrow 1$.\label{step:1/2S}
%						
%					}
					
				}
				
				
				
			}
			
		}
		
		
		
		
	}
	Return the (primal) solution $x$ and the (dual) solution $y$. 
\end{algorithm}

In time $t$, the algorithm considers a {\em violating set} of pages $S \subseteq \cP$. This set violates the primal constraint in \eqref{eq:LP2} corresponding to $t$ and $S$. %and is of maximum cardinality amongst all such sets.   %We remove all pages in $Q$ except $p_t$ from cache and 
We increase the variable $y_{t}(S)$ continuously and at the same time increase variables $x_p(r(p,t))$ for all pages in $S$ (except for $p_t$) that are not fully evicted yet. The increase rate is a function of $Y_p(r(p,t))$, where %recall that $Y_p(r(p,t))$
\begin{equation}
	\label{eq:Y2}
	Y_p(j) =\sum_{t \in I(p,j)~} \sum_{S \subseteq \cP \big|p \in S-p_t} y_t(S)
\end{equation}  is the left hand side of the corresponding dual constraint of $p$ and $j = r(p,t)$ in \eqref{eq:LP2} (analogously to \eqref{eq:Y}). This growth function has an exponential dependence on $Y_p(r(p,t))$, scaled by the cost of page $c(p)$, the number of pages, and the total cover demand. The growth of variable $x_p(j)$ stops once it reaches $1$. %$\frac{1}{2}$.   



%To compute $x_p(j)$, we continuously increase the (dual) variable $y_t(Q)$, corresponding to the non-valid set of pages in the cache at time $t$. At the same time, we increase the probability $x_p(t)$ (and equivalently, $X_p(t)$) according to an exponential function that depends on $y_p(j)$, where recall that $y_p(j)$ is the sum over all dual variables corresponding to $p$ in the interval $R(p,j)$. Once a valid distribution is constructed, we sample a page according to $x(t)$ and achieving a valid cache state once again. We also remove from the cache pages $p'$ for which $X_{p'}(t) \geq 1$. The pseudocode of the algorithm is given in \Cref{alg:fractional}. 





Once the primal constraint corresponding to $t,S$ is satisfied, there are two cases. If $x$ becomes feasible, the algorithm proceeds to the next time step. Otherwise, the algorithm repeats the above process with a new set $S'$ for which $x$ does not satisfy its constraint. The pseudocode of the algorithm is given in \Cref{alg:fractionalS}. 


\subsubsection*{Analysis of \Cref{alg:fractionalS}}
In the following we analyze the competitive performance of the algorithm.  We first prove the feasibility of the primal and dual solutions $x,y$. %Note that the algorithm returns a feasible primal solution $x$ by \Cref{step:InWhileS}. 


	\begin{lemma}
	\label{lem:PFeasibleS}
	\Cref{alg:fractionalS} returns a feasible primal solution $x$ to \eqref{eq:LPS}. 
\end{lemma}

\begin{proof}
	Let $t \in T$ and $S \in \cS(t)$. We show that $x$ satisfies the corresponding primal constraint of $t,S$. Let $U = \left\{p \in S~|~x_p(r(p,t)) = 1\right\}$ be the set of fully-evicted pages in $S$ (observe that $p_t \notin U$ since $x_{p_t}(r(p,t)) = 0$). We use the following auxiliary claims. 
	
		\begin{claim}
		\label{claim:r<44}
		%$\cS\left(S \setminus U\right) \subseteq \left\{S' \in \cS(S)~|~S' \setminus U\right\}$
		For every $S' \in \cC(S \setminus U)$ it holds that $S' \cup U \in \cC(S)$.
	\end{claim}
	\begin{claimproof}
		It holds that 
		$$f \left(S \setminus \left(S' \cup U\right)\right) = f \left((S \setminus U) \setminus S'\right) \leq k.$$
		The inequality holds since $S' \in \cC(S \setminus U)$ (see the definition of the sets $\cC(S),\cC(S \setminus U)$ above \eqref{eq:q(S)}). By the definition of $\cC(S)$ it holds that $S' \cup U \in \cC(S)$.
	\end{claimproof}
	The following claim relies on Claim~\ref{claim:r<44}. 
		\begin{claim}
		\label{claim:r<45}
		%$\cS\left(S \setminus U\right) \subseteq \left\{S' \in \cS(S)~|~S' \setminus U\right\}$
	$q(S) \leq |U|+q(S \setminus U)$. 
	\end{claim}
	\begin{claimproof}
		It holds that 
		\begin{equation*}
			\begin{aligned}
			q(S) ={} & \min_{S' \in \cC(S)} |S'| 
			\\={} & \min_{S' \in \cC(S)} \left( |S' \cap U|+|S' \setminus U| \right) 
			\\\leq{} &  
			\min_{S' \in \cC(S) \text{ s.t. } U \subseteq S'} \left( |S' \cap U|+|S' \setminus U| \right) 
				\\={} &  
				|U|+\min_{S' \in \cC(S) \text{ s.t. } U \subseteq S'}|S' \setminus U| 
					\\\leq{} &  
				|U|+\min_{S' \in \cC(S \setminus U)} |S'| 
					\\={} &  
				|U|+q(S \setminus U). 
			\end{aligned}
		\end{equation*} The first inequality holds by restricting the minimum to a subset of the set $\cC(S)$; note that there is $S' \in \cC(S)$ such that $U \subseteq S'$ since $S \in \cC(S)$ and $U \subseteq S$. The second inequality follows from Claim~\ref{claim:r<44}. 
	\end{claimproof}
	
	Therefore, 
	$$\sum_{p \in S-p_t} x_p(r(p,t)) = \sum_{p \in U} x_p(r(p,t))+\sum_{p \in (S\setminus U)-p_t} x_p(r(p,t)) \geq |U|+q(S \setminus U) \geq q(S).$$
	The first inequality holds by the following. First, since every page in $U$ is fully evicted. Second, since $$\sum_{p \in (S\setminus U)-p_t} x_p(r(p,t)) \geq q(S \setminus U)$$ as non of the pages in $S \setminus U$ is fully-evicted; thus, the constraint of $t,S-U$ must be satisfied by the stopping condition of \Cref{step:InWhileS}.  The last inequality follows from Claim~\ref{claim:r<45}. 
\end{proof}


We now show the feasibility of the dual solution $y$ using the growth rate of the variables in the algorithm. 
	\begin{lemma}
	\label{lem:FFeasibleS}
	\Cref{alg:fractionalS} returns a feasible dual solution to \eqref{eq:LPS}. 
\end{lemma}

\begin{proof}
	%To prove that the dual $y$ is feasible, 
	Consider some $p \in \cP$ and $j \in [n_p]$. By the update rate of $x_p(j)$ in \Cref{step:fxS} it follows that 
	\begin{equation}
		\label{eq:FromSAPS}
		\frac{1}{\mu} \cdot \left(     \exp \left(  \frac{\ln (\mu+1)}{c(p)} \cdot Y_p(j)\right)-1\right) = x_p(j) \leq 1.
	\end{equation} 
	The inequality holds since $Y_p(j)$ does not increase in time $t$ if $p$ is already fully evicted at this time by \Cref{step:InWhileS}. Thus, by simplifying the expression in \eqref{eq:FromSAPS} it holds that 
	$$\exp \left( \frac{\ln (\mu+1)}{c(p)} \cdot Y_p(j) \right) \leq \mu+1.$$
	by taking $\ln$ over both sides it follows that $Y_p(j) \leq c(p)$. Hence, $y$ is a feasible dual solution. 
\end{proof}

We finally prove the competitive ratio of the algorithm.

%	\begin{lemma}
%	\label{thm:logL1S}
%	%	\Cref{alg:fractional} is $O(\log (n))$-\textnormal{competitive}. 
%	The cost of $x$ is bounded by $O(\log (\mu))$ times the value of the dual $y$.  
%\end{lemma}

\subsubsection*{Proof of \Cref{thm:MU}}

\begin{proof}
	Consider an infinitesimal increase in the value of the dual solution $y$. Specifically, assume that the algorithm chooses a set $S$ in \Cref{step:fQS} for time time step $t$, and that the dual variable $y_t(S)$ increases infinitesimally by $d y_t(S)$. Let $dx$ and $dy$ denote the infinitesimal change in the objective value of $x$ and $y$, respectively. We bound the increase $dx$ in $x$ as a function of the increase $dy$.   
	\begin{equation}
		\label{eq:dxS}
		\begin{aligned}
			dx ={} & \sum_{p \in S-p_t} d x_p(r(p,t)) \cdot c(p) \\
			={} & \sum_{p \in S-p_t}  \frac{d x_p(r(p,t)) \cdot c(p) \cdot d y_t(S)}{d y_t(S)} \\
			={} & \sum_{p \in S-p_t} \ln \left(\mu+1 \right)\ \cdot \left( x_p(r(p,t))+\frac{1}{\mu} \right) \cdot d y_t(S). 
		\end{aligned}
	\end{equation} The first equality holds since the increase in $y_t(S)$ induces an increase only for the primal variables corresponding to pages in $S-p_t$. The last equality follows from the growth rate of a variable $x_p(r(p,t))$, for some $p \in S-p_t$ as a result of the growth in $y_t(S)$. %(for more details see~\cite{bansal2012primal,bansal2012randomized}). 
	We separately analyze the two expressions in \eqref{eq:dxS}. First, since at the considered moment in time $y_t(S)$ increases, by \Cref{step:InWhileS} it holds that 
	\begin{equation}
		\label{eq:InwhileFS}
		\sum_{p \in S-p_t} x_p(r(p,t)) \leq q(S)%\cdot g_Q(p) \leq 	\sum_{p \in \cP-p_t} x_p(r(p,t)) \cdot g_Q(p) < N-g(Q)
		%\sum_{p \in Q-p_t} x_p(r(p,t)) < 1.
	\end{equation} 
	
	For the second expression, let $S' \in \cC(S)$ such that $|S'| = q(S)$; clearly, by \eqref{eq:q(S)} there is such $S'$. By the definition of $\cC(S)$ it holds that $f(S \setminus S') \leq k$; thus, $|S \setminus S'| \leq \mu$. Therefore,  
	\begin{equation}
		\label{eq:WidthInS}
		\sum_{p \in S-p_t} \frac{1}{\mu} = \frac{ \left| S \right|-1}{\mu} = \frac{ \left| S' \right|+\left|S \setminus S'\right|-1}{\mu} \leq \frac{ q(S)+\mu-1}{\mu} \leq q(S)+1 \leq 2 \cdot q(S).  
	\end{equation} For the last inequality, note that $q(S) \geq 1$ since we assume that $y_t(S)$ increases implying that the corresponding constraint of $t$ and $S$ is not satisfied. Therefore, by \eqref{eq:dxS}, \eqref{eq:InwhileFS}, and~\eqref{eq:WidthInS},
	\begin{equation}
		\label{eq:PDESS}
		\begin{aligned}
			dx \leq{} &  \ln \left(\mu+1 \right) \cdot d y_t(S) \cdot  \left(	\sum_{p \in S-p_t} x_p(r(p,t)) + 	\sum_{p \in S-p_t} \frac{1}{\mu} \right) \\
			={} & 3 \cdot \ln \left(\mu+1 \right) \cdot d y_t(S) \cdot q(S)\\
			={} & O(\log (\mu)) \cdot dy.
		\end{aligned}
	\end{equation} Thus, by \eqref{eq:PDESS}, every increase in $y$ increases $x$ by a factor of at most $O(\log (\mu))$. By \Cref{lem:FFeasibleS,lem:PFeasibleS} the algorithm also returns feasible primal and dual solutions, implying the statement of the theorem.  
	%Finally, note that if $x_p(j) \geq \frac{1}{2}$ then we immediately increase $x_p(j)$ to $1$; this increase the total cost of $x$ by a factor of $2$ w.r.t. the value of $y$. 
	%Since $x$ and $y$ are feasible primal and dual solutions by \Cref{lem:FFeasible}, the proof follows from \Cref{lem:relaxation2}. 
\end{proof}